% !TeX program = lualatex
\documentclass[12pt,aspectratio=169]{beamer}
\usepackage[utf8]{inputenc}
\usepackage[T1]{fontenc}
\usepackage{amsmath,amsfonts,amssymb}
\usepackage{graphicx}
\usepackage{xcolor}
\usepackage{hyperref}
\usepackage{anyfontsize}
\usepackage{lmodern}

% ====== EXTREME FONT REDUCTION ======
\renewcommand{\tiny}{\fontsize{4}{5}\selectfont}
\renewcommand{\scriptsize}{\fontsize{5}{6}\selectfont}
\renewcommand{\footnotesize}{\fontsize{6}{7}\selectfont}
\renewcommand{\small}{\fontsize{7}{8}\selectfont}
\renewcommand{\normalsize}{\fontsize{8}{9}\selectfont}
\renewcommand{\large}{\fontsize{9}{10}\selectfont}
\renewcommand{\Large}{\fontsize{10}{11}\selectfont}
\renewcommand{\LARGE}{\fontsize{11}{13}\selectfont}
\renewcommand{\huge}{\fontsize{12}{14}\selectfont}
\renewcommand{\Huge}{\fontsize{14}{17}\selectfont}

% ====== COLOR THEME ======
\definecolor{redcorrect}{RGB}{200,30,30}
\definecolor{bluecorrect}{RGB}{30,80,200}
\definecolor{greencheck}{RGB}{0,150,50}
\definecolor{lightgray}{RGB}{245,245,245}
\definecolor{darkgray}{RGB}{50,50,50}
\definecolor{softyellow}{RGB}{255,248,200}
\definecolor{steppurple}{RGB}{150,50,200}
\definecolor{deepblue}{RGB}{20,60,150}
\definecolor{darkred}{RGB}{150,20,20}
\definecolor{orangewarning}{RGB}{255,140,0}

\setbeamercolor{title}{fg=white,bg=redcorrect}
\setbeamercolor{frametitle}{fg=white,bg=redcorrect}
\setbeamercolor{block title}{fg=white,bg=bluecorrect}
\setbeamercolor{block body}{fg=darkgray,bg=lightgray}
\setbeamercolor{normal text}{fg=darkgray}
\usecolortheme{default}

% ====== CUSTOM COMMANDS ======
\newcommand{\correct}[1]{\textcolor{greencheck}{\textbf{\checkmark}} #1}
\newcommand{\keypoint}[1]{\textcolor{deepblue}{\textbf{#1}}}
\newcommand{\hardconcept}[1]{\textcolor{darkred}{\textbf{#1}}}
\newcommand{\tipbox}[1]{\fcolorbox{bluecorrect}{softyellow}{\parbox{0.88\textwidth}{\centering #1}}}
\newcommand{\stepbox}[1]{%
	\begin{center}
		\fcolorbox{darkgray}{white}{%
			\parbox{0.90\textwidth}{\vspace{0.05cm} #1 \vspace{0.05cm}}%
		}
	\end{center}
}
\newcommand{\wrongbox}[1]{\fcolorbox{redcorrect}{red!10}{\parbox{0.88\textwidth}{\centering \textcolor{redcorrect}{\textbf{\times}} #1}}}
\newcommand{\highlight}[1]{\textcolor{redcorrect}{\textbf{#1}}}
\newcommand{\bluetext}[1]{\textcolor{bluecorrect}{\textbf{#1}}}

\begin{document}
	
	% ================================================================
	% SLIDE 1: TITLE
	% ================================================================
	\begin{frame}
		\centering
		\vspace{0.8cm}
		{\fontsize{16}{19}\selectfont \textbf{Tools and Technologies to Fight Financial Crime}}\\[0.15cm]
		{\fontsize{12}{14}\selectfont Domain D — Complete Q\&A Review}\\[0.3cm]
		{\fontsize{9}{11}\selectfont Understanding the \textit{Regulatory Principles} Behind Each Answer}\\[0.3cm]
		{\fontsize{8}{10}\selectfont Compliance Training Department}
		\vspace{0.2cm}
		\rule{4cm}{0.5pt}
	\end{frame}
	
	% ================================================================
	% SLIDE 2: AGENDA
	% ================================================================
	\begin{frame}
		\frametitle{\fontsize{11}{13}\selectfont Agenda — Domain D Topics}
		\scriptsize
		\begin{enumerate}
			\item Tools Across Customer Lifecycle (Q1-Q4)
			\item Data Quality, Integrity, Access, Definitions (Q5-Q8)
			\item Customer Experience and Friction (Q9-Q12)
			\item Onboarding — Digital KYC, E-KYC, Biometrics (Q13-Q16)
			\item Onboarding — External Data Sources (Q17-Q20)
			\item Onboarding — Name/Sanctions Screening (Q21-Q24)
			\item Ongoing — Periodic Refresh and Perpetual KYC (Q25-Q28)
			\item Ongoing — Adverse Media Screening (Q29-Q32)
			\item Ongoing — Batch Screening and Fuzzy Logic (Q33-Q36)
			\item Ongoing — Transaction/Payment Screening (Q37-Q40)
			\item Ongoing — Rules-Based TM and Scenario Tuning (Q41-Q44)
			\item Ongoing — Transitioning to AI/ML Tools (Q45-Q48)
			\item Ongoing — Investigations and Automation (Q49-Q52)
			\item Ongoing — Network Analysis Tools (Q53-Q56)
			\item Privacy and Data Protection Technologies (Q57-Q60)
			\item Choosing and Integrating AFC Tools (Q61-Q64)
			\item Operational Effectiveness and Resource Prioritization (Q65-Q68)
		\end{enumerate}
	\end{frame}
	
	% ================================================================
	% SLIDE 3-6: TOOLS ACROSS CUSTOMER LIFECYCLE
	% ================================================================
	\begin{frame}
		\frametitle{\fontsize{11}{13}\selectfont Q1: Tools Across Customer Lifecycle}
		\begin{block}{\fontsize{8}{10}\selectfont Topic: Tools Across Customer Lifecycle}
			\scriptsize
			What tools are available to implement controls across the customer lifecycle?
		\end{block}
		\vspace{0.05cm}
		\stepbox{\scriptsize
			\keypoint{Correct Answer:} Onboarding tools (KYC, screening), ongoing monitoring tools (transaction monitoring, periodic reviews), and exit controls (offboarding, account closure).
		}
		\vspace{0.05cm}
		\scriptsize
		\hardconcept{Core Concept:} The customer lifecycle has three phases:
		\par \textbf{Onboarding}: Digital KYC, E-KYC, biometrics, sanctions screening.
		\par \textbf{Ongoing}: Transaction monitoring, periodic reviews, perpetual KYC.
		\par \textbf{Exit}: Offboarding controls, account closure, FIU liaison.
		\vspace{0.03cm}
		\wrongbox{\scriptsize \textbf{Common Trap:} "AML tools are only needed at onboarding" — INCORRECT. Controls must be applied \textbf{throughout} the lifecycle.}
		\vspace{0.03cm}
		\tipbox{\scriptsize \textbf{Key Insight:} The lifecycle approach ensures \textbf{continuous} compliance.}
	\end{frame}
	
	\begin{frame}
		\frametitle{\fontsize{11}{13}\selectfont Q2: Onboarding Tools}
		\begin{block}{\fontsize{8}{10}\selectfont Topic: Onboarding Tools}
			\scriptsize
			What tools are used at customer onboarding for AFC compliance?
		\end{block}
		\vspace{0.05cm}
		\stepbox{\scriptsize
			\keypoint{Correct Answer:} Digital identity verification, biometric authentication, sanctions screening, PEP screening, and adverse media checks.
		}
		\vspace{0.03cm}
		\tipbox{\scriptsize \textbf{Key Insight:} Onboarding is the \textbf{critical first line} of defense.}
	\end{frame}
	
	\begin{frame}
		\frametitle{\fontsize{11}{13}\selectfont Q3: Ongoing Monitoring Tools}
		\begin{block}{\fontsize{8}{10}\selectfont Topic: Ongoing Monitoring Tools}
			\scriptsize
			What tools are used for ongoing AFC monitoring?
		\end{block}
		\vspace{0.05cm}
		\stepbox{\scriptsize
			\keypoint{Correct Answer:} Transaction monitoring systems, periodic KYC refresh, perpetual KYC, and adverse media screening.
		}
		\vspace{0.03cm}
		\tipbox{\scriptsize \textbf{Key Insight:} Ongoing monitoring detects \textbf{evolving} risks.}
	\end{frame}
	
	\begin{frame}
		\frametitle{\fontsize{11}{13}\selectfont Q4: Exit Controls}
		\begin{block}{\fontsize{8}{10}\selectfont Topic: Exit Controls}
			\scriptsize
			What controls apply when exiting a customer relationship?
		\end{block}
		\vspace{0.05cm}
		\stepbox{\scriptsize
			\keypoint{Correct Answer:} Offboarding procedures, account closure protocols, and liaison with FIUs/law enforcement if needed.
		}
		\vspace{0.03cm}
		\tipbox{\scriptsize \textbf{Key Insight:} Exits must be \textbf{managed} and \textbf{documented}.}
	\end{frame}
	
	% ================================================================
	% SLIDE 7-10: DATA QUALITY
	% ================================================================
	\begin{frame}
		\frametitle{\fontsize{11}{13}\selectfont Q5: Data Quality — Importance}
		\begin{block}{\fontsize{8}{10}\selectfont Topic: Data Quality, Integrity, Access}
			\scriptsize
			Why is data quality important in AFC compliance?
		\end{block}
		\vspace{0.05cm}
		\stepbox{\scriptsize
			\keypoint{Correct Answer:} Poor data quality leads to false positives, missed detections, and regulatory penalties.
		}
		\vspace{0.05cm}
		\scriptsize
		\hardconcept{Core Concept:} \textbf{GIGO} — Garbage In, Garbage Out. Data quality dimensions:
		\par \textbf{Accuracy}: Data correctly represents reality.
		\par \textbf{Completeness}: All required data is present.
		\par \textbf{Consistency}: Data is consistent across systems.
		\par \textbf{Timeliness}: Data is current and up-to-date.
		\vspace{0.03cm}
		\wrongbox{\scriptsize \textbf{Common Trap:} "Automated systems compensate for poor data quality" — INCORRECT. Systems are \textbf{only as good} as their data inputs.}
		\vspace{0.03cm}
		\tipbox{\scriptsize \textbf{Key Insight:} Under \textbf{BCBS 239}, banks must have effective data quality controls.}
	\end{frame}
	
	\begin{frame}
		\frametitle{\fontsize{11}{13}\selectfont Q6: Data Integrity}
		\begin{block}{\fontsize{8}{10}\selectfont Topic: Data Integrity}
			\scriptsize
			What is data integrity in the AFC context?
		\end{block}
		\vspace{0.05cm}
		\stepbox{\scriptsize
			\keypoint{Correct Answer:} Ensuring that data is accurate, complete, and has not been altered or corrupted.
		}
		\vspace{0.03cm}
		\tipbox{\scriptsize \textbf{Key Insight:} Data integrity is essential for \textbf{regulatory reporting}.}
	\end{frame}
	
	\begin{frame}
		\frametitle{\fontsize{11}{13}\selectfont Q7: Data Access Controls}
		\begin{block}{\fontsize{8}{10}\selectfont Topic: Data Access}
			\scriptsize
			What data access controls are needed in AFC?
		\end{block}
		\vspace{0.05cm}
		\stepbox{\scriptsize
			\keypoint{Correct Answer:} Role-based access controls, segregation of duties, and audit trails.
		}
		\vspace{0.03cm}
		\tipbox{\scriptsize \textbf{Key Insight:} Access controls prevent \textbf{unauthorized} data use.}
	\end{frame}
	
	\begin{frame}
		\frametitle{\fontsize{11}{13}\selectfont Q8: Data Definitions and Taxonomy}
		\begin{block}{\fontsize{8}{10}\selectfont Topic: Data Definitions}
			\scriptsize
			Why is a common data taxonomy important in AFC?
		\end{block}
		\vspace{0.05cm}
		\stepbox{\scriptsize
			\keypoint{Correct Answer:} A common taxonomy ensures consistent data interpretation across systems and teams.
		}
		\vspace{0.03cm}
		\tipbox{\scriptsize \textbf{Key Insight:} A data dictionary provides \textbf{consistent definitions}.}
	\end{frame}
	
	% ================================================================
	% SLIDE 11-14: CUSTOMER EXPERIENCE
	% ================================================================
	\begin{frame}
		\frametitle{\fontsize{11}{13}\selectfont Q9: Customer Experience and Friction}
		\begin{block}{\fontsize{8}{10}\selectfont Topic: Customer Experience}
			\scriptsize
			How can institutions balance compliance with customer experience?
		\end{block}
		\vspace{0.05cm}
		\stepbox{\scriptsize
			\keypoint{Correct Answer:} By applying a risk-based approach with proportionate friction — lower friction for low-risk customers, higher for high-risk.
		}
		\vspace{0.05cm}
		\scriptsize
		\hardconcept{Core Concept:}
		\par \textbf{Low-Risk Customers}: Simplified onboarding, minimal friction.
		\par \textbf{Standard Customers}: Standard CDD, reasonable friction.
		\par \textbf{High-Risk Customers}: Enhanced due diligence, higher friction.
		\vspace{0.03cm}
		\wrongbox{\scriptsize \textbf{Common Trap:} "All customers should have the same onboarding experience" — INCORRECT. A risk-based approach requires \textbf{differentiated} experiences.}
		\vspace{0.03cm}
		\tipbox{\scriptsize \textbf{Key Insight:} \textbf{Risk-based friction} balances compliance and experience.}
	\end{frame}
	
	\begin{frame}
		\frametitle{\fontsize{11}{13}\selectfont Q10: Digital Onboarding Experience}
		\begin{block}{\fontsize{8}{10}\selectfont Topic: Digital Onboarding}
			\scriptsize
			How can digital onboarding improve customer experience?
		\end{block}
		\vspace{0.05cm}
		\stepbox{\scriptsize
			\keypoint{Correct Answer:} By enabling remote, fast, and convenient onboarding with minimal manual intervention.
		}
		\vspace{0.03cm}
		\tipbox{\scriptsize \textbf{Key Insight:} Digital onboarding is \textbf{expected} by modern customers.}
	\end{frame}
	
	\begin{frame}
		\frametitle{\fontsize{11}{13}\selectfont Q11: Friction in High-Risk Scenarios}
		\begin{block}{\fontsize{8}{10}\selectfont Topic: Friction}
			\scriptsize
			When should institutions increase friction for customers?
		\end{block}
		\vspace{0.05cm}
		\stepbox{\scriptsize
			\keypoint{Correct Answer:} For high-risk customers, complex transactions, or when red flags are identified.
		}
		\vspace{0.03cm}
		\tipbox{\scriptsize \textbf{Key Insight:} Friction is a \textbf{control} for high-risk situations.}
	\end{frame}
	
	\begin{frame}
		\frametitle{\fontsize{11}{13}\selectfont Q12: Risk-Based Authentication}
		\begin{block}{\fontsize{8}{10}\selectfont Topic: Risk-Based Authentication}
			\scriptsize
			What is risk-based authentication?
		\end{block}
		\vspace{0.05cm}
		\stepbox{\scriptsize
			\keypoint{Correct Answer:} Authentication that prompts additional checks only when user behavior deviates from expectations.
		}
		\vspace{0.03cm}
		\tipbox{\scriptsize \textbf{Key Insight:} RBA \textbf{optimizes} user experience while maintaining security.}
	\end{frame}
	
	% ================================================================
	% SLIDE 15-18: ONBOARDING — DIGITAL KYC
	% ================================================================
	\begin{frame}
		\frametitle{\fontsize{11}{13}\selectfont Q13: Digital Onboarding Solutions}
		\begin{block}{\fontsize{8}{10}\selectfont Topic: Digital Onboarding}
			\scriptsize
			What are the latest digital onboarding solutions for KYC?
		\end{block}
		\vspace{0.05cm}
		\stepbox{\scriptsize
			\keypoint{Correct Answer:} E-KYC, digital identity verification, facial recognition, liveness detection, biometrics, and geolocation.
		}
		\vspace{0.05cm}
		\scriptsize
		\hardconcept{Core Concept:}
		\par \textbf{E-KYC}: Electronic identity verification.
		\par \textbf{Digital Identity}: Verified digital credentials.
		\par \textbf{Facial Recognition}: Comparing selfie to ID photo.
		\par \textbf{Liveness Detection}: Ensuring the person is physically present.
		\par \textbf{Biometrics}: Fingerprint, voice, or facial recognition.
		\par \textbf{Geolocation}: Verifying location matches profile.
		\vspace{0.03cm}
		\tipbox{\scriptsize \textbf{Key Insight:} Multi-factor authentication is \textbf{essential} for robust identity verification.}
	\end{frame}
	
	\begin{frame}
		\frametitle{\fontsize{11}{13}\selectfont Q14: E-KYC}
		\begin{block}{\fontsize{8}{10}\selectfont Topic: E-KYC}
			\scriptsize
			What is E-KYC?
		\end{block}
		\vspace{0.05cm}
		\stepbox{\scriptsize
			\keypoint{Correct Answer:} Electronic Know Your Customer — digital identity verification using electronic documents and identity databases.
		}
		\vspace{0.03cm}
		\tipbox{\scriptsize \textbf{Key Insight:} E-KYC enables \textbf{remote} onboarding.}
	\end{frame}
	
	\begin{frame}
		\frametitle{\fontsize{11}{13}\selectfont Q15: Facial Recognition and Liveness}
		\begin{block}{\fontsize{8}{10}\selectfont Topic: Facial Recognition}
			\scriptsize
			How does liveness detection work in facial recognition?
		\end{block}
		\vspace{0.05cm}
		\stepbox{\scriptsize
			\keypoint{Correct Answer:} It analyzes subtle facial movements to confirm the person is physically present and not using a photo or video.
		}
		\vspace{0.03cm}
		\tipbox{\scriptsize \textbf{Key Insight:} Liveness detection prevents \textbf{spoofing} attacks.}
	\end{frame}
	
	\begin{frame}
		\frametitle{\fontsize{11}{13}\selectfont Q16: Biometrics in Onboarding}
		\begin{block}{\fontsize{8}{10}\selectfont Topic: Biometrics}
			\scriptsize
			What biometric methods are used in customer onboarding?
		\end{block}
		\vspace{0.05cm}
		\stepbox{\scriptsize
			\keypoint{Correct Answer:} Fingerprint scanning, facial recognition, voice recognition, and behavioral biometrics.
		}
		\vspace{0.03cm}
		\wrongbox{\scriptsize \textbf{Common Trap:} "Biometrics are completely secure and cannot be compromised" — INCORRECT. Biometrics can be \textbf{stolen} and cannot be reset.}
		\vspace{0.03cm}
		\tipbox{\scriptsize \textbf{Key Insight:} Biometric data must be \textbf{protected} like any other personal data.}
	\end{frame}
	
	% ================================================================
	% SLIDE 19-22: ONBOARDING — EXTERNAL DATA SOURCES
	% ================================================================
	\begin{frame}
		\frametitle{\fontsize{11}{13}\selectfont Q17: External Data Sources for KYC}
		\begin{block}{\fontsize{8}{10}\selectfont Topic: External Data Sources}
			\scriptsize
			What external data sources are used for KYC and screening?
		\end{block}
		\vspace{0.05cm}
		\stepbox{\scriptsize
			\keypoint{Correct Answer:} Credit reference agencies, beneficial ownership registers, adverse media, criminal records, and government identity checks.
		}
		\vspace{0.05cm}
		\scriptsize
		\hardconcept{Core Concept:}
		\par \textbf{Credit Reference Agencies}: Verify identity and financial history.
		\par \textbf{Beneficial Ownership Registers}: Identify ultimate owners.
		\par \textbf{Adverse Media}: Search for negative news.
		\par \textbf{Criminal Records}: Check for criminal history.
		\par \textbf{Government Identity Checks}: Verify against official records.
		\vspace{0.03cm}
		\tipbox{\scriptsize \textbf{Key Insight:} External sources provide \textbf{third-party validation} of customer information.}
	\end{frame}
	
	\begin{frame}
		\frametitle{\fontsize{11}{13}\selectfont Q18: Automated vs Manual Data Sources}
		\begin{block}{\fontsize{8}{10}\selectfont Topic: Automated vs Manual}
			\scriptsize
			What is the difference between automated and manual external data sources?
		\end{block}
		\vspace{0.05cm}
		\stepbox{\scriptsize
			\keypoint{Correct Answer:} Automated sources are queried in real-time via APIs; manual sources require human intervention to access.
		}
		\vspace{0.03cm}
		\tipbox{\scriptsize \textbf{Key Insight:} Automation improves \textbf{efficiency} and \textbf{speed}.}
	\end{frame}
	
	\begin{frame}
		\frametitle{\fontsize{11}{13}\selectfont Q19: Beneficial Ownership Registers}
		\begin{block}{\fontsize{8}{10}\selectfont Topic: BO Registers}
			\scriptsize
			How are beneficial ownership registers used in KYC?
		\end{block}
		\vspace{0.05cm}
		\stepbox{\scriptsize
			\keypoint{Correct Answer:} To identify the ultimate beneficial owners of legal entities and verify their identities.
		}
		\vspace{0.03cm}
		\tipbox{\scriptsize \textbf{Key Insight:} BO registers support \textbf{transparency} and \textbf{CDD}.}
	\end{frame}
	
	\begin{frame}
		\frametitle{\fontsize{11}{13}\selectfont Q20: Government Identity Checks}
		\begin{block}{\fontsize{8}{10}\selectfont Topic: Government Checks}
			\scriptsize
			What are government identity checks used for in KYC?
		\end{block}
		\vspace{0.05cm}
		\stepbox{\scriptsize
			\keypoint{Correct Answer:} To verify the authenticity of identity documents against official government records.
		}
		\vspace{0.03cm}
		\tipbox{\scriptsize \textbf{Key Insight:} Government checks provide \textbf{authoritative} verification.}
	\end{frame}
	
	% ================================================================
	% SLIDE 23-26: ONBOARDING — SANCTIONS SCREENING
	% ================================================================
	\begin{frame}
		\frametitle{\fontsize{11}{13}\selectfont Q21: Name and Customer Screening}
		\begin{block}{\fontsize{8}{10}\selectfont Topic: Name Screening}
			\scriptsize
			What is name screening in the AFC context?
		\end{block}
		\vspace{0.05cm}
		\stepbox{\scriptsize
			\keypoint{Correct Answer:} The process of checking customer names against sanctions lists, watchlists, and other risk databases.
		}
		\vspace{0.05cm}
		\scriptsize
		\hardconcept{Core Concept:} Screening sources:
		\par \textbf{UN Sanctions Lists}: Global sanctions.
		\par \textbf{OFAC SDN List}: US sanctions.
		\par \textbf{EU Sanctions Lists}: EU sanctions.
		\par \textbf{PEP Lists}: Politically exposed persons.
		\par \textbf{Internal Watchlists}: Institution-specific.
		\vspace{0.03cm}
		\tipbox{\scriptsize \textbf{Key Insight:} Screening should be \textbf{real-time} for onboarding.}
	\end{frame}
	
	\begin{frame}
		\frametitle{\fontsize{11}{13}\selectfont Q22: Sanctions Lists Used}
		\begin{block}{\fontsize{8}{10}\selectfont Topic: Sanctions Lists}
			\scriptsize
			What sanctions lists are commonly used for customer screening?
		\end{block}
		\vspace{0.05cm}
		\stepbox{\scriptsize
			\keypoint{Correct Answer:} UN sanctions lists, OFAC SDN List, EU sanctions lists, and other internal and external watchlists.
		}
		\vspace{0.03cm}
		\tipbox{\scriptsize \textbf{Key Insight:} Institutions must screen against \textbf{all applicable} lists.}
	\end{frame}
	
	\begin{frame}
		\frametitle{\fontsize{11}{13}\selectfont Q23: Internal and External Watchlists}
		\begin{block}{\fontsize{8}{10}\selectfont Topic: Watchlists}
			\scriptsize
			What is the difference between internal and external watchlists?
		\end{block}
		\vspace{0.05cm}
		\stepbox{\scriptsize
			\keypoint{Correct Answer:} External watchlists are from regulators/governments; internal watchlists are created by the institution based on previous investigations.
		}
		\vspace{0.03cm}
		\tipbox{\scriptsize \textbf{Key Insight:} Internal watchlists capture \textbf{institution-specific} risks.}
	\end{frame}
	
	\begin{frame}
		\frametitle{\fontsize{11}{13}\selectfont Q24: Fraud and TF Watchlists}
		\begin{block}{\fontsize{8}{10}\selectfont Topic: Fraud/TF Lists}
			\scriptsize
			What are fraud and terrorist financing watchlists?
		\end{block}
		\vspace{0.05cm}
		\stepbox{\scriptsize
			\keypoint{Correct Answer:} Lists identifying known fraudsters, terrorist financiers, and associated entities that should be screened against.
		}
		\vspace{0.03cm}
		\tipbox{\scriptsize \textbf{Key Insight:} These lists are \textbf{essential} for fraud and TF detection.}
	\end{frame}
	
	% ================================================================
	% SLIDE 27-30: ONGOING — PERPETUAL KYC
	% ================================================================
	\begin{frame}
		\frametitle{\fontsize{11}{13}\selectfont Q25: Periodic Refresh and Perpetual KYC}
		\begin{block}{\fontsize{8}{10}\selectfont Topic: Perpetual KYC}
			\scriptsize
			What is the difference between periodic refresh and perpetual KYC?
		\end{block}
		\vspace{0.05cm}
		\stepbox{\scriptsize
			\keypoint{Correct Answer:} Periodic refresh is conducted at fixed intervals; perpetual KYC is triggered by real-time data changes and events.
		}
		\vspace{0.05cm}
		\scriptsize
		\hardconcept{Core Concept:}
		\par \textbf{Periodic Refresh}: Annual or biennial reviews.
		\par \textbf{Perpetual KYC}: Continuous monitoring with trigger-based reviews.
		\par \textbf{Prioritization}: High-risk customers reviewed more frequently.
		\par \textbf{AI/ML}: Automates data monitoring and event detection.
		\vspace{0.03cm}
		\wrongbox{\scriptsize \textbf{Common Trap:} "Perpetual KYC replaces all periodic reviews" — INCORRECT. Perpetual KYC \textbf{complements} periodic reviews.}
		\vspace{0.03cm}
		\tipbox{\scriptsize \textbf{Key Insight:} Perpetual KYC enables \textbf{real-time} risk detection.}
	\end{frame}
	
	\begin{frame}
		\frametitle{\fontsize{11}{13}\selectfont Q26: Prioritization in KYC Refresh}
		\begin{block}{\fontsize{8}{10}\selectfont Topic: Prioritization}
			\scriptsize
			How should KYC refresh be prioritized?
		\end{block}
		\vspace{0.05cm}
		\stepbox{\scriptsize
			\keypoint{Correct Answer:} By applying a risk-based approach — higher-risk customers refreshed more frequently than lower-risk customers.
		}
		\vspace{0.03cm}
		\tipbox{\scriptsize \textbf{Key Insight:} Prioritization ensures resources are \textbf{allocated efficiently}.}
	\end{frame}
	
	\begin{frame}
		\frametitle{\fontsize{11}{13}\selectfont Q27: AI/ML in Perpetual KYC}
		\begin{block}{\fontsize{8}{10}\selectfont Topic: AI/ML in PKYC}
			\scriptsize
			How does AI/ML support perpetual KYC?
		\end{block}
		\vspace{0.05cm}
		\stepbox{\scriptsize
			\keypoint{Correct Answer:} By automatically monitoring for data changes, identifying risk triggers, and prioritizing reviews.
		}
		\vspace{0.03cm}
		\tipbox{\scriptsize \textbf{Key Insight:} AI/ML enables \textbf{scalable} perpetual KYC.}
	\end{frame}
	
	\begin{frame}
		\frametitle{\fontsize{11}{13}\selectfont Q28: Trigger-Based Reviews}
		\begin{block}{\fontsize{8}{10}\selectfont Topic: Trigger-Based Reviews}
			\scriptsize
			What triggers a customer review in perpetual KYC?
		\end{block}
		\vspace{0.05cm}
		\stepbox{\scriptsize
			\keypoint{Correct Answer:} Significant changes in behavior, new risk indicators, regulatory updates, or system alerts.
		}
		\vspace{0.03cm}
		\tipbox{\scriptsize \textbf{Key Insight:} Triggers are \textbf{event-driven}, not schedule-driven.}
	\end{frame}
	
	% ================================================================
	% SLIDE 31-34: ONGOING — ADVERSE MEDIA
	% ================================================================
	\begin{frame}
		\frametitle{\fontsize{11}{13}\selectfont Q29: Adverse Media Screening Tools}
		\begin{block}{\fontsize{8}{10}\selectfont Topic: Adverse Media}
			\scriptsize
			What are adverse media screening tools?
		\end{block}
		\vspace{0.05cm}
		\stepbox{\scriptsize
			\keypoint{Correct Answer:} Tools that search for negative news about customers, including criminal activity, corruption, and regulatory issues.
		}
		\vspace{0.05cm}
		\scriptsize
		\hardconcept{Core Concept:} Adverse media sources:
		\par \textbf{Global News}: Reuters, Bloomberg, BBC.
		\par \textbf{Regulatory Databases}: Enforcement actions.
		\par \textbf{Court Records}: Legal proceedings.
		\par \textbf{Specialized Databases}: Risk intelligence providers.
		\vspace{0.03cm}
		\tipbox{\scriptsize \textbf{Key Insight:} Adverse media is a \textbf{critical} component of CDD and EDD.}
	\end{frame}
	
	\begin{frame}
		\frametitle{\fontsize{11}{13}\selectfont Q30: AI/ML in Adverse Media}
		\begin{block}{\fontsize{8}{10}\selectfont Topic: AI/ML Adverse Media}
			\scriptsize
			How does AI/ML improve adverse media screening?
		\end{block}
		\vspace{0.05cm}
		\stepbox{\scriptsize
			\keypoint{Correct Answer:} By automating search across multiple sources, identifying relevant articles, and reducing false positives.
		}
		\vspace{0.03cm}
		\tipbox{\scriptsize \textbf{Key Insight:} AI/ML enables \textbf{comprehensive} and \textbf{efficient} adverse media screening.}
	\end{frame}
	
	\begin{frame}
		\frametitle{\fontsize{11}{13}\selectfont Q31: Adverse Media Red Flags}
		\begin{block}{\fontsize{8}{10}\selectfont Topic: Red Flags}
			\scriptsize
			What are red flags in adverse media screening?
		\end{block}
		\vspace{0.05cm}
		\stepbox{\scriptsize
			\keypoint{Correct Answer:} Articles about fraud, corruption, sanctions violations, regulatory fines, or criminal charges.
		}
		\vspace{0.03cm}
		\tipbox{\scriptsize \textbf{Key Insight:} Red flags trigger \textbf{EDD} or \textbf{investigation}.}
	\end{frame}
	
	\begin{frame}
		\frametitle{\fontsize{11}{13}\selectfont Q32: Continuous Adverse Media Monitoring}
		\begin{block}{\fontsize{8}{10}\selectfont Topic: Continuous Monitoring}
			\scriptsize
			Why is continuous adverse media monitoring important?
		\end{block}
		\vspace{0.05cm}
		\stepbox{\scriptsize
			\keypoint{Correct Answer:} Customer risk profiles can change quickly based on new adverse information.
		}
		\vspace{0.03cm}
		\tipbox{\scriptsize \textbf{Key Insight:} \textbf{Continuous} monitoring detects emerging risks.}
	\end{frame}
	
	% ================================================================
	% SLIDE 35-38: ONGOING — BATCH SCREENING
	% ================================================================
	\begin{frame}
		\frametitle{\fontsize{11}{13}\selectfont Q33: Batch Screening for Sanctions}
		\begin{block}{\fontsize{8}{10}\selectfont Topic: Batch Screening}
			\scriptsize
			What is batch screening for sanctions?
		\end{block}
		\vspace{0.05cm}
		\stepbox{\scriptsize
			\keypoint{Correct Answer:} Processing the entire customer database against sanctions lists in bulk, typically on a periodic basis.
		}
		\vspace{0.05cm}
		\scriptsize
		\hardconcept{Core Concept:}
		\par \textbf{Batch Screening}: Periodic, volume-oriented screening.
		\par \textbf{Real-Time Screening}: Individual screening at point of interaction.
		\par \textbf{Hybrid Approach}: Both batch and real-time screening.
		\vspace{0.03cm}
		\wrongbox{\scriptsize \textbf{Common Trap:} "Batch screening replaces real-time screening" — INCORRECT. Both are \textbf{complementary}.}
		\vspace{0.03cm}
		\tipbox{\scriptsize \textbf{Key Insight:} Batch screening detects \textbf{changes} in customer risk status.}
	\end{frame}
	
	\begin{frame}
		\frametitle{\fontsize{11}{13}\selectfont Q34: Fuzzy Logic in Screening}
		\begin{block}{\fontsize{8}{10}\selectfont Topic: Fuzzy Logic}
			\scriptsize
			What is fuzzy logic in sanctions screening?
		\end{block}
		\vspace{0.05cm}
		\stepbox{\scriptsize
			\keypoint{Correct Answer:} Matching logic that accounts for name variations, spelling differences, and typos to identify potential matches.
		}
		\vspace{0.03cm}
		\tipbox{\scriptsize \textbf{Key Insight:} Fuzzy logic reduces \textbf{false negatives} from name variations.}
	\end{frame}
	
	\begin{frame}
		\frametitle{\fontsize{11}{13}\selectfont Q35: AI/ML in Sanctions Screening}
		\begin{block}{\fontsize{8}{10}\selectfont Topic: AI/ML Screening}
			\scriptsize
			How does AI/ML improve sanctions screening?
		\end{block}
		\vspace{0.05cm}
		\stepbox{\scriptsize
			\keypoint{Correct Answer:} By reducing false positives, improving match accuracy, and learning from investigator decisions.
		}
		\vspace{0.03cm}
		\tipbox{\scriptsize \textbf{Key Insight:} AI/ML \textbf{improves} screening over time.}
	\end{frame}
	
	\begin{frame}
		\frametitle{\fontsize{11}{13}\selectfont Q36: List Management}
		\begin{block}{\fontsize{8}{10}\selectfont Topic: List Management}
			\scriptsize
			What is list management in sanctions screening?
		\end{block}
		\vspace{0.05cm}
		\stepbox{\scriptsize
			\keypoint{Correct Answer:} The process of maintaining, updating, and managing sanctions lists and watchlists used for screening.
		}
		\vspace{0.03cm}
		\tipbox{\scriptsize \textbf{Key Insight:} Lists must be \textbf{current} and \textbf{accurate}.}
	\end{frame}
	
	% ================================================================
	% SLIDE 39-42: ONGOING — TRANSACTION SCREENING
	% ================================================================
	\begin{frame}
		\frametitle{\fontsize{11}{13}\selectfont Q37: Transaction/Payment Screening}
		\begin{block}{\fontsize{8}{10}\selectfont Topic: Transaction Screening}
			\scriptsize
			What is transaction/payment screening?
		\end{block}
		\vspace{0.05cm}
		\stepbox{\scriptsize
			\keypoint{Correct Answer:} Real-time screening of payment messages against sanctions lists, watchlists, and other risk databases.
		}
		\vspace{0.05cm}
		\scriptsize
		\hardconcept{Core Concept:} Payment message types:
		\par \textbf{SWIFT}: Financial messaging network.
		\par \textbf{Blockchain Transactions}: Cryptoasset transfers.
		\par \textbf{Payment Systems}: ACH, SEPA, Fedwire.
		\vspace{0.03cm}
		\tipbox{\scriptsize \textbf{Key Insight:} Transaction screening is \textbf{critical} for sanctions compliance.}
	\end{frame}
	
	\begin{frame}
		\frametitle{\fontsize{11}{13}\selectfont Q38: SWIFT Payment Screening}
		\begin{block}{\fontsize{8}{10}\selectfont Topic: SWIFT}
			\scriptsize
			How does SWIFT payment screening work?
		\end{block}
		\vspace{0.05cm}
		\stepbox{\scriptsize
			\keypoint{Correct Answer:} SWIFT messages are screened for sanctioned parties, restricted entities, and suspicious patterns.
		}
		\vspace{0.03cm}
		\tipbox{\scriptsize \textbf{Key Insight:} SWIFT messages contain \textbf{structured data} for screening.}
	\end{frame}
	
	\begin{frame}
		\frametitle{\fontsize{11}{13}\selectfont Q39: Blockchain Transaction Screening}
		\begin{block}{\fontsize{8}{10}\selectfont Topic: Blockchain Screening}
			\scriptsize
			What are the challenges in screening blockchain transactions?
		\end{block}
		\vspace{0.05cm}
		\stepbox{\scriptsize
			\keypoint{Correct Answer:} Pseudonymity, use of privacy coins, mixers, and rapid movement across addresses.
		}
		\vspace{0.03cm}
		\tipbox{\scriptsize \textbf{Key Insight:} Blockchain analytics tools address these \textbf{challenges}.}
	\end{frame}
	
	\begin{frame}
		\frametitle{\fontsize{11}{13}\selectfont Q40: AI/ML in Transaction Screening}
		\begin{block}{\fontsize{8}{10}\selectfont Topic: AI/ML Transaction Screening}
			\scriptsize
			How does AI/ML improve transaction screening?
		\end{block}
		\vspace{0.05cm}
		\stepbox{\scriptsize
			\keypoint{Correct Answer:} By detecting patterns, reducing false positives, and identifying emerging typologies.
		}
		\vspace{0.03cm}
		\tipbox{\scriptsize \textbf{Key Insight:} AI/ML \textbf{adapts} to evolving criminal methods.}
	\end{frame}
	
	% ================================================================
	% SLIDE 43-46: ONGOING — RULES-BASED TM
	% ================================================================
	\begin{frame}
		\frametitle{\fontsize{11}{13}\selectfont Q41: Rules-Based Transaction Monitoring}
		\begin{block}{\fontsize{8}{10}\selectfont Topic: Rules-Based TM}
			\scriptsize
			What is rules-based transaction monitoring?
		\end{block}
		\vspace{0.05cm}
		\stepbox{\scriptsize
			\keypoint{Correct Answer:} Monitoring that uses predefined rules and thresholds to identify potentially suspicious transactions.
		}
		\vspace{0.05cm}
		\scriptsize
		\hardconcept{Core Concept:} TM components:
		\par \textbf{Customer Segmentation}: Grouping by risk profile.
		\par \textbf{Scenario Coverage}: Rules for different risk types.
		\par \textbf{Threshold Setting}: Trigger points for alerts.
		\par \textbf{Statistical Testing}: Above The Line (ATL) / Below The Line (BTL).
		\vspace{0.03cm}
		\wrongbox{\scriptsize \textbf{Common Trap:} "More alerts mean better compliance" — INCORRECT. \textbf{Quality} of alerts matters more than quantity.}
		\vspace{0.03cm}
		\tipbox{\scriptsize \textbf{Key Insight:} Rules-based TM is the \textbf{foundation} of monitoring programs.}
	\end{frame}
	
	\begin{frame}
		\frametitle{\fontsize{11}{13}\selectfont Q42: Customer Segmentation}
		\begin{block}{\fontsize{8}{10}\selectfont Topic: Segmentation}
			\scriptsize
			Why is customer segmentation important in TM?
		\end{block}
		\vspace{0.05cm}
		\stepbox{\scriptsize
			\keypoint{Correct Answer:} Different customer segments have different risk profiles and require different monitoring thresholds.
		}
		\vspace{0.03cm}
		\tipbox{\scriptsize \textbf{Key Insight:} Segmentation enables \textbf{proportionate} monitoring.}
	\end{frame}
	
	\begin{frame}
		\frametitle{\fontsize{11}{13}\selectfont Q43: Threshold Setting and Tuning}
		\begin{block}{\fontsize{8}{10}\selectfont Topic: Threshold Setting}
			\scriptsize
			How are TM thresholds set and tuned?
		\end{block}
		\vspace{0.05cm}
		\stepbox{\scriptsize
			\keypoint{Correct Answer:} Through statistical analysis, testing (ATL/BTL), and regular tuning based on alert productivity metrics.
		}
		\vspace{0.03cm}
		\tipbox{\scriptsize \textbf{Key Insight:} Tuning is an \textbf{iterative} process.}
	\end{frame}
	
	\begin{frame}
		\frametitle{\fontsize{11}{13}\selectfont Q44: Model Risk Management}
		\begin{block}{\fontsize{8}{10}\selectfont Topic: Model Risk}
			\scriptsize
			What is model risk management in AFC?
		\end{block}
		\vspace{0.05cm}
		\stepbox{\scriptsize
			\keypoint{Correct Answer:} Governance and controls over models used for AML, including monitoring, validation, and documentation.
		}
		\vspace{0.03cm}
		\tipbox{\scriptsize \textbf{Key Insight:} Under \textbf{SR 11-7}, banks must have model risk management.}
	\end{frame}
	
	% ================================================================
	% SLIDE 47-50: ONGOING — TRANSITIONING TO AI/ML
	% ================================================================
	\begin{frame}
		\frametitle{\fontsize{11}{13}\selectfont Q45: Transitioning to AI/ML Tools}
		\begin{block}{\fontsize{8}{10}\selectfont Topic: AI/ML Transition}
			\scriptsize
			How can institutions transition from rules-based to AI/ML-based monitoring?
		\end{block}
		\vspace{0.05cm}
		\stepbox{\scriptsize
			\keypoint{Correct Answer:} Gradual transition starting with hybrid approaches, then moving to primary and secondary AI/ML tools.
		}
		\vspace{0.05cm}
		\scriptsize
		\hardconcept{Core Concept:} Transition stages:
		\par \textbf{Stage 1}: Rules-based primary, AI/ML secondary.
		\par \textbf{Stage 2}: AI/ML primary, rules-based secondary.
		\par \textbf{Stage 3}: Fully AI/ML-driven monitoring.
		\vspace{0.03cm}
		\tipbox{\scriptsize \textbf{Key Insight:} Transition should be \textbf{phased} and \textbf{validated}.}
	\end{frame}
	
	\begin{frame}
		\frametitle{\fontsize{11}{13}\selectfont Q46: Primary AI/ML Tools}
		\begin{block}{\fontsize{8}{10}\selectfont Topic: Primary AI/ML}
			\scriptsize
			What are primary AI/ML tools in AFC?
		\end{block}
		\vspace{0.05cm}
		\stepbox{\scriptsize
			\keypoint{Correct Answer:} Machine learning models that detect patterns and anomalies without predefined rules.
		}
		\vspace{0.03cm}
		\tipbox{\scriptsize \textbf{Key Insight:} Primary AI/ML \textbf{learns} from data.}
	\end{frame}
	
	\begin{frame}
		\frametitle{\fontsize{11}{13}\selectfont Q47: Secondary AI/ML Tools}
		\begin{block}{\fontsize{8}{10}\selectfont Topic: Secondary AI/ML}
			\scriptsize
			What are secondary AI/ML tools in AFC?
		\end{block}
		\vspace{0.05cm}
		\stepbox{\scriptsize
			\keypoint{Correct Answer:} ML models that enhance rules-based systems by reducing false positives and prioritizing alerts.
		}
		\vspace{0.03cm}
		\tipbox{\scriptsize \textbf{Key Insight:} Secondary AI/ML \textbf{augments} existing systems.}
	\end{frame}
	
	\begin{frame}
		\frametitle{\fontsize{11}{13}\selectfont Q48: Transition Challenges}
		\begin{block}{\fontsize{8}{10}\selectfont Topic: Transition Challenges}
			\scriptsize
			What are challenges in transitioning to AI/ML tools?
		\end{block}
		\vspace{0.05cm}
		\stepbox{\scriptsize
			\keypoint{Correct Answer:} Data quality, model validation, explainability, and regulatory acceptance.
		}
		\vspace{0.03cm}
		\tipbox{\scriptsize \textbf{Key Insight:} \textbf{Explainability} is critical for regulatory acceptance.}
	\end{frame}
	
	% ================================================================
	% SLIDE 51-54: ONGOING — INVESTIGATIONS
	% ================================================================
	\begin{frame}
		\frametitle{\fontsize{11}{13}\selectfont Q49: Investigations — Open Source Tools}
		\begin{block}{\fontsize{8}{10}\selectfont Topic: Investigations}
			\scriptsize
			What open-source tools are used in AFC investigations?
		\end{block}
		\vspace{0.05cm}
		\stepbox{\scriptsize
			\keypoint{Correct Answer:} Public records searches, social media monitoring, company registers, and OSINT tools.
		}
		\vspace{0.05cm}
		\scriptsize
		\hardconcept{Core Concept:} OSINT sources:
		\par \textbf{Public Records}: Court documents, property records.
		\par \textbf{Social Media}: LinkedIn, Facebook, Twitter.
		\par \textbf{Company Registers}: Corporate ownership information.
		\par \textbf{News}: Current and historical articles.
		\vspace{0.03cm}
		\tipbox{\scriptsize \textbf{Key Insight:} OSINT provides \textbf{valuable context} for investigations.}
	\end{frame}
	
	\begin{frame}
		\frametitle{\fontsize{11}{13}\selectfont Q50: Automation in Investigations}
		\begin{block}{\fontsize{8}{10}\selectfont Topic: Automation}
			\scriptsize
			How does automation support AFC investigators?
		\end{block}
		\vspace{0.05cm}
		\stepbox{\scriptsize
			\keypoint{Correct Answer:} By automating data gathering, analysis, and workflow management to improve efficiency.
		}
		\vspace{0.03cm}
		\tipbox{\scriptsize \textbf{Key Insight:} Automation allows investigators to focus on \textbf{high-value} work.}
	\end{frame}
	
	\begin{frame}
		\frametitle{\fontsize{11}{13}\selectfont Q51: AI/ML in Investigations}
		\begin{block}{\fontsize{8}{10}\selectfont Topic: AI/ML in Investigations}
			\scriptsize
			How does AI/ML support AFC investigations?
		\end{block}
		\vspace{0.05cm}
		\stepbox{\scriptsize
			\keypoint{Correct Answer:} By identifying patterns, connecting data points, and prioritizing leads.
		}
		\vspace{0.03cm}
		\tipbox{\scriptsize \textbf{Key Insight:} AI/ML \textbf{augments} human investigators.}
	\end{frame}
	
	\begin{frame}
		\frametitle{\fontsize{11}{13}\selectfont Q52: Investigation Workflow Tools}
		\begin{block}{\fontsize{8}{10}\selectfont Topic: Workflow Tools}
			\scriptsize
			What workflow tools support AFC investigations?
		\end{block}
		\vspace{0.05cm}
		\stepbox{\scriptsize
			\keypoint{Correct Answer:} Case management systems, workflow automation, and collaboration tools.
		}
		\vspace{0.03cm}
		\tipbox{\scriptsize \textbf{Key Insight:} Workflow tools ensure \textbf{consistency} and \textbf{completeness}.}
	\end{frame}
	
	% ================================================================
	% SLIDE 55-58: ONGOING — NETWORK ANALYSIS
	% ================================================================
	\begin{frame}
		\frametitle{\fontsize{11}{13}\selectfont Q53: Network Analysis Tools}
		\begin{block}{\fontsize{8}{10}\selectfont Topic: Network Analysis}
			\scriptsize
			What are network analysis tools in AFC?
		\end{block}
		\vspace{0.05cm}
		\stepbox{\scriptsize
			\keypoint{Correct Answer:} Tools that visualize and analyze connections between entities to identify criminal networks.
		}
		\vspace{0.05cm}
		\scriptsize
		\hardconcept{Core Concept:} Network analysis:
		\par \textbf{Node Analysis}: Identifying key entities.
		\par \textbf{Link Analysis}: Identifying connections.
		\par \textbf{Clustering}: Identifying groups.
		\par \textbf{Visualization}: Visual representation of networks.
		\vspace{0.03cm}
		\tipbox{\scriptsize \textbf{Key Insight:} Network analysis reveals \textbf{hidden connections}.}
	\end{frame}
	
	\begin{frame}
		\frametitle{\fontsize{11}{13}\selectfont Q54: Identifying Criminal Links}
		\begin{block}{\fontsize{8}{10}\selectfont Topic: Criminal Links}
			\scriptsize
			How do network analysis tools identify links to criminality?
		\end{block}
		\vspace{0.05cm}
		\stepbox{\scriptsize
			\keypoint{Correct Answer:} By connecting entities through shared addresses, transactions, phone numbers, and other relationships.
		}
		\vspace{0.03cm}
		\tipbox{\scriptsize \textbf{Key Insight:} Connections reveal \textbf{suspicious relationships}.}
	\end{frame}
	
	\begin{frame}
		\frametitle{\fontsize{11}{13}\selectfont Q55: AI/ML in Network Analysis}
		\begin{block}{\fontsize{8}{10}\selectfont Topic: AI/ML Network Analysis}
			\scriptsize
			How does AI/ML improve network analysis?
		\end{block}
		\vspace{0.05cm}
		\stepbox{\scriptsize
			\keypoint{Correct Answer:} By detecting hidden patterns and relationships that human analysts might miss.
		}
		\vspace{0.03cm}
		\tipbox{\scriptsize \textbf{Key Insight:} AI/ML \textbf{augments} human pattern recognition.}
	\end{frame}
	
	\begin{frame}
		\frametitle{\fontsize{11}{13}\selectfont Q56: Clustering in Network Analysis}
		\begin{block}{\fontsize{8}{10}\selectfont Topic: Clustering}
			\scriptsize
			What is clustering in network analysis?
		\end{block}
		\vspace{0.05cm}
		\stepbox{\scriptsize
			\keypoint{Correct Answer:} Grouping entities that share common characteristics, indicating potential criminal networks.
		}
		\vspace{0.03cm}
		\tipbox{\scriptsize \textbf{Key Insight:} Clustering reveals \textbf{shared patterns}.}
	\end{frame}
	
	% ================================================================
	% SLIDE 59-62: PRIVACY AND DATA PROTECTION
	% ================================================================
	\begin{frame}
		\frametitle{\fontsize{11}{13}\selectfont Q57: Privacy and Data Protection Technologies}
		\begin{block}{\fontsize{8}{10}\selectfont Topic: Privacy Technologies}
			\scriptsize
			What technologies help navigate privacy and data protection rules in AFC?
		\end{block}
		\vspace{0.05cm}
		\stepbox{\scriptsize
			\keypoint{Correct Answer:} Data anonymization, pseudonymization, differential privacy, and secure multi-party computation.
		}
		\vspace{0.05cm}
		\scriptsize
		\hardconcept{Core Concept:}
		\par \textbf{Anonymization}: Removing identifiers.
		\par \textbf{Pseudonymization}: Replacing identifiers with tokens.
		\par \textbf{Differential Privacy}: Adding noise to prevent re-identification.
		\par \textbf{Secure MPC}: Computing without revealing data.
		\vspace{0.03cm}
		\wrongbox{\scriptsize \textbf{Common Trap:} "Privacy technologies are not needed for AML compliance" — INCORRECT. AML data processing \textbf{must} comply with privacy laws.}
		\vspace{0.03cm}
		\tipbox{\scriptsize \textbf{Key Insight:} Privacy technologies enable \textbf{data protection} while maintaining effectiveness.}
	\end{frame}
	
	\begin{frame}
		\frametitle{\fontsize{11}{13}\selectfont Q58: GDPR-Compliant AFC Tools}
		\begin{block}{\fontsize{8}{10}\selectfont Topic: GDPR Compliance}
			\scriptsize
			What features should GDPR-compliant AFC tools have?
		\end{block}
		\vspace{0.05cm}
		\stepbox{\scriptsize
			\keypoint{Correct Answer:} Data minimization, purpose limitation, access controls, and audit trails.
		}
		\vspace{0.03cm}
		\tipbox{\scriptsize \textbf{Key Insight:} GDPR applies to \textbf{all} personal data processing.}
	\end{frame}
	
	\begin{frame}
		\frametitle{\fontsize{11}{13}\selectfont Q59: Data Minimization in Tools}
		\begin{block}{\fontsize{8}{10}\selectfont Topic: Data Minimization}
			\scriptsize
			How is data minimization applied in AFC tools?
		\end{block}
		\vspace{0.05cm}
		\stepbox{\scriptsize
			\keypoint{Correct Answer:} Collecting and processing only the data necessary for the specific AFC purpose.
		}
		\vspace{0.03cm}
		\tipbox{\scriptsize \textbf{Key Insight:} Minimization reduces \textbf{privacy risk}.}
	\end{frame}
	
	\begin{frame}
		\frametitle{\fontsize{11}{13}\selectfont Q60: Cross-Border Data Transfers}
		\begin{block}{\fontsize{8}{10}\selectfont Topic: Cross-Border Transfers}
			\scriptsize
			How do AFC tools handle cross-border data transfers?
		\end{block}
		\vspace{0.05cm}
		\stepbox{\scriptsize
			\keypoint{Correct Answer:} Using appropriate safeguards such as Standard Contractual Clauses or adequacy decisions.
		}
		\vspace{0.03cm}
		\tipbox{\scriptsize \textbf{Key Insight:} Transfers must be \textbf{legally compliant}.}
	\end{frame}
	
	% ================================================================
	% SLIDE 63-66: CHOOSING AND INTEGRATING TOOLS
	% ================================================================
	\begin{frame}
		\frametitle{\fontsize{11}{13}\selectfont Q61: Choosing the Best AFC Tool}
		\begin{block}{\fontsize{8}{10}\selectfont Topic: Choosing Tools}
			\scriptsize
			What should organizations consider when choosing an AFC tool?
		\end{block}
		\vspace{0.05cm}
		\stepbox{\scriptsize
			\keypoint{Correct Answer:} Risk profile, institutional size, existing systems, budget, regulatory requirements, and scalability.
		}
		\vspace{0.05cm}
		\scriptsize
		\hardconcept{Core Concept:} Selection criteria:
		\par \textbf{Functionality}: Does it meet requirements?
		\par \textbf{Integration}: Can it connect to existing systems?
		\par \textbf{Scalability}: Will it grow with the institution?
		\par \textbf{Cost}: Is it cost-effective?
		\par \textbf{Vendor Support}: Is support available?
		\vspace{0.03cm}
		\wrongbox{\scriptsize \textbf{Common Trap:} "The most expensive tool is always the best" — INCORRECT. The \textbf{right} tool depends on the \textbf{institution's needs}.}
		\vspace{0.03cm}
		\tipbox{\scriptsize \textbf{Key Insight:} Tool selection should be \textbf{needs-driven}.}
	\end{frame}
	
	\begin{frame}
		\frametitle{\fontsize{11}{13}\selectfont Q62: Integration Considerations}
		\begin{block}{\fontsize{8}{10}\selectfont Topic: Integration}
			\scriptsize
			What should be considered when integrating AFC tools?
		\end{block}
		\vspace{0.05cm}
		\stepbox{\scriptsize
			\keypoint{Correct Answer:} Data mapping, system compatibility, data quality, and change management.
		}
		\vspace{0.03cm}
		\tipbox{\scriptsize \textbf{Key Insight:} Integration is a \textbf{critical success factor}.}
	\end{frame}
	
	\begin{frame}
		\frametitle{\fontsize{11}{13}\selectfont Q63: Vendor Assessment}
		\begin{block}{\fontsize{8}{10}\selectfont Topic: Vendor Assessment}
			\scriptsize
			What should be assessed when evaluating AFC tool vendors?
		\end{block}
		\vspace{0.05cm}
		\stepbox{\scriptsize
			\keypoint{Correct Answer:} Vendor reputation, financial stability, security posture, and compliance track record.
		}
		\vspace{0.03cm}
		\tipbox{\scriptsize \textbf{Key Insight:} Vendor assessment \textbf{mitigates} third-party risk.}
	\end{frame}
	
	\begin{frame}
		\frametitle{\fontsize{11}{13}\selectfont Q64: Proof of Concept}
		\begin{block}{\fontsize{8}{10}\selectfont Topic: Proof of Concept}
			\scriptsize
			Why should organizations conduct a proof of concept (POC) for AFC tools?
		\end{block}
		\vspace{0.05cm}
		\stepbox{\scriptsize
			\keypoint{Correct Answer:} To validate the tool's effectiveness, integration, and fit with institutional requirements.
		}
		\vspace{0.03cm}
		\tipbox{\scriptsize \textbf{Key Insight:} A POC \textbf{mitigates} implementation risk.}
	\end{frame}
	
	% ================================================================
	% SLIDE 67-70: OPERATIONAL EFFECTIVENESS
	% ================================================================
	\begin{frame}
		\frametitle{\fontsize{11}{13}\selectfont Q65: Operational Effectiveness}
		\begin{block}{\fontsize{8}{10}\selectfont Topic: Operational Effectiveness}
			\scriptsize
			What is operational effectiveness in AFC compliance?
		\end{block}
		\vspace{0.05cm}
		\stepbox{\scriptsize
			\keypoint{Correct Answer:} Ensuring AFC controls are efficient and achieve their intended purpose without waste.
		}
		\vspace{0.05cm}
		\scriptsize
		\hardconcept{Core Concept:} Effectiveness dimensions:
		\par \textbf{Detection}: Identifying suspicious activity.
		\par \textbf{Efficiency}: Minimizing false positives.
		\par \textbf{Cost-Effectiveness}: Achieving results within budget.
		\par \textbf{Resource Utilization}: Using resources optimally.
		\vspace{0.03cm}
		\tipbox{\scriptsize \textbf{Key Insight:} Effective controls are \textbf{efficient} controls.}
	\end{frame}
	
	\begin{frame}
		\frametitle{\fontsize{11}{13}\selectfont Q66: Resource Prioritization}
		\begin{block}{\fontsize{8}{10}\selectfont Topic: Resource Prioritization}
			\scriptsize
			How should resources be prioritized in AFC compliance?
		\end{block}
		\vspace{0.05cm}
		\stepbox{\scriptsize
			\keypoint{Correct Answer:} By applying a risk-based approach that allocates resources to the highest-risk areas.
		}
		\vspace{0.03cm}
		\tipbox{\scriptsize \textbf{Key Insight:} \textbf{High-risk} areas get more resources.}
	\end{frame}
	
	\begin{frame}
		\frametitle{\fontsize{11}{13}\selectfont Q67: Investment in AFC Tools}
		\begin{block}{\fontsize{8}{10}\selectfont Topic: Investment}
			\scriptsize
			How should institutions prioritize investment in AFC tools?
		\end{block}
		\vspace{0.05cm}
		\stepbox{\scriptsize
			\keypoint{Correct Answer:} Based on risk assessment, regulatory requirements, and cost-benefit analysis.
		}
		\vspace{0.03cm}
		\tipbox{\scriptsize \textbf{Key Insight:} Investment should be \textbf{risk-driven}.}
	\end{frame}
	
	\begin{frame}
		\frametitle{\fontsize{11}{13}\selectfont Q68: Measuring Effectiveness}
		\begin{block}{\fontsize{8}{10}\selectfont Topic: Measuring Effectiveness}
			\scriptsize
			How is the effectiveness of AFC tools measured?
		\end{block}
		\vspace{0.05cm}
		\stepbox{\scriptsize
			\keypoint{Correct Answer:} Through KPIs, KRIs, alert productivity metrics, and regulatory feedback.
		}
		\vspace{0.03cm}
		\tipbox{\scriptsize \textbf{Key Insight:} Measurement enables \textbf{continuous improvement}.}
	\end{frame}
	
	% ================================================================
	% SLIDE 71-74: SUMMARY
	% ================================================================
	\begin{frame}
		\frametitle{\fontsize{11}{13}\selectfont SUMMARY: Key Principles — Domain D}
		\scriptsize
		\begin{block}{\fontsize{9}{11}\selectfont Tools and Technologies to Fight Financial Crime:}
			\vspace{0.05cm}
			\scriptsize
			\begin{tabular}{|p{0.35\textwidth}|p{0.60\textwidth}|}
				\hline
				\keypoint{Concept} & \keypoint{Application} \\
				\hline
				Lifecycle Tools & Onboarding, ongoing, exit controls \\
				\hline
				Data Quality & Accuracy, completeness, consistency, timeliness \\
				\hline
				Customer Experience & Risk-based friction; RBA \\
				\hline
				Digital Onboarding & E-KYC, facial recognition, liveness, biometrics \\
				\hline
				External Data Sources & Credit agencies, BO registers, adverse media \\
				\hline
				Name/Sanctions Screening & UN, OFAC, EU, watchlists, fuzzy logic \\
				\hline
				Perpetual KYC & Real-time, trigger-based, risk-prioritized \\
				\hline
				Adverse Media & OSINT, AI/ML, continuous monitoring \\
				\hline
				Batch Screening & Periodic, volume-oriented, list management \\
				\hline
				Transaction Screening & SWIFT, blockchain, AI/ML \\
				\hline
				Rules-Based TM & Segmentation, thresholds, ATL/BTL, model risk \\
				\hline
				AI/ML Transition & Hybrid → primary → fully AI-driven \\
				\hline
				Investigations & OSINT, automation, workflow tools \\
				\hline
				Network Analysis & Link analysis, clustering, visualization \\
				\hline
				Privacy Technologies & Anonymization, pseudonymization, differential privacy \\
				\hline
				Tool Selection & Needs-driven, integration, POC \\
				\hline
				Operational Effectiveness & Risk-based, efficient, measured \\
				\hline
			\end{tabular}
			\vspace{0.05cm}
		\end{block}
		\vspace{0.05cm}
		\tipbox{\scriptsize \textbf{FINAL LESSON:} Domain D tests understanding of \textbf{how technology supports AFC} — from onboarding to ongoing monitoring to investigations. Understand the \textit{tools} and the \textit{answers} become obvious.}
	\end{frame}
	
	\begin{frame}
		\frametitle{\fontsize{11}{13}\selectfont EXAM TIPS — Domain D}
		\scriptsize
		\begin{block}{\fontsize{9}{11}\selectfont Common Traps to Avoid:}
			\vspace{0.05cm}
			\scriptsize
			\begin{itemize}
				\item \highlight{Data Quality}: GIGO — systems are only as good as their data.
				\item \highlight{Perpetual KYC}: \textbf{Complements} periodic reviews, does not replace them.
				\item \highlight{Batch vs Real-Time}: \textbf{Complementary}, not substitutes.
				\item \highlight{More Alerts}: More alerts ≠ better compliance — \textbf{quality} matters.
				\item \highlight{Fuzzy Logic}: Reduces \textbf{false negatives}, not false positives.
				\item \highlight{AI/ML Transition}: Gradual, \textbf{phased}, and \textbf{validated}.
				\item \highlight{Privacy}: AML does not exempt institutions from \textbf{GDPR}.
				\item \highlight{Tool Selection}: \textbf{Needs-driven}, not vendor-driven.
				\item \highlight{Network Analysis}: Reveals \textbf{hidden connections} between entities.
				\item \highlight{Model Risk}: Under \textbf{SR 11-7}, models must be validated.
			\end{itemize}
			\vspace{0.05cm}
		\end{block}
		\vspace{0.05cm}
		\tipbox{\scriptsize \textbf{Key Insight:} Focus on understanding the \textit{role} of each technology and \textit{when} it should be applied.}
	\end{frame}
	
	\begin{frame}
		\centering
		\vspace{2.5cm}
		{\fontsize{16}{19}\selectfont \textbf{Thank You}} \\[0.2cm]
		{\fontsize{11}{13}\selectfont Keep learning. Keep growing.} \\[0.15cm]
		\rule{3cm}{0.5pt} \\[0.15cm]
		{\fontsize{8}{10}\selectfont \textit{Technology is a powerful enabler of AFC compliance when applied correctly.}}
		\vspace{0.3cm}
		{\fontsize{7}{9}\selectfont \textcolor{darkgray}{End of Domain D Review — 75 Slides}}
	\end{frame}
	
\end{document}